\documentclass[10pt]{article}

\usepackage[utf8]{inputenc}

\usepackage[T1]{fontenc}

\usepackage{amsmath}

\usepackage{amsfonts}

\usepackage{amssymb}

\usepackage[version=4]{mhchem}

\usepackage{stmaryrd}

\usepackage{bbold}



\begin{document}

бых $\lambda \in R$ и $a \in A$. При $A=B$ получается определение а н д омо ю и з м а. Все эндоморфизмы полигона $A$ образуют моноид, и $A$ можно рассматривать как II. над ним.



Эквивалентность $\theta$ на $R$-полигоне $A$ наз. к о нг р у ә ц ц и й, если $(a, b) \in \theta$ влечет $(\lambda a, \lambda b) \in \theta$ для любого $\lambda \in R$. Множество классов конгруэнции $\theta$ естественным образом превращается в $R$-полигон, называемый фа к то р по ли г о і о м полигона $A$ и обозначаемый через $A / \theta$. Если $A-$ полигон над $R$, то на $R$ можно определить отношение Ann $A$, положив $(\lambda, \mu) \in \operatorname{Ann} A$, если $\lambda a=\mu a$ для всех $a \in A$. Отношение $\operatorname{Ann} A$ оказывается конгруэнцией моноида $R$, и $A$ естественным образом превращается в П. над фактормоноидом $R / \operatorname{Ann} A$. Если полигон $A$ возник из нек-рого автомата, то описанный переход равносилен «склеиванию" одинаковым образом действующих последовательностей входных сигналов. Наряду с обычными для универсальных алгебр конструкциями прямого и подпрямого произведения, в теории П. рассматривается важная для алгебраич. теории автоматов конструкция сплетения. Свободное произведение (или копроизведение) П. совпадает с их дизъюнктным объединением.



На П. можно смотреть как на неаддитивный аналог модуля над кольцом, что служит богатым источником задач теории П. В частности, установлена связь П. с радикалами в полугруппах и исследуются связи свойств моноида со свойствами П. над ним. Напр., все левыө $R$-полигоны проективны тогда и только тогда, когда $R$ - одноэлементная группа, а инъективность всех П. над коммутативным моноидом $R$ равносильна наличию в $R$ нуля и порождаемости всех его идеалов идемпотентами (ср. Гомологическая классификация колец).



Если $R$ - моноид с нулем 0 , то можно говорить об $R$-полигоне $A$ с нулем как об $R$-полигоне с отмеченной точкой $u$, причем $0 a=u$ для всех $a \in A$. Теория П. с нулем имеет нек-рые особенности.



Каждый П. можно рассматривать как функтор из однообъектной категории в категорию множеств.



Лum.: [1] Алгебраическая теория автоматов, языков и полугрупп, пер. с англ., М., 1975; [2] К л и ф ф о р д А., П р е c-

то н Г., Алгебраическая теория полугрупп, пер. с англ., т. 2 , то н Г., Алгебраическая теория полугрупп, пер. с англ., т. 2 ,

М., $1972:$ [3] С к о р н я к о в Л. А., в сб.: Модули, в. 3, Новосиб., 1973 , с. $22-27 ;$ [4] Итоги науки и техники. Алгебра. Топология. Геометрия, т. 14, М., 1976, с. 57-190.



ПОЛИКРУГ, по ли ц и ли н д р,-область



\[

\begin{gathered}

\Delta=\Delta\left(a=\left(a_{1}, \ldots, a_{n}\right), r=\left(r_{1}, \ldots, r_{n}\right)=\right. \\

=\left\{z=\left(z_{1}, \ldots, z_{n}\right) \in \mathbb{C}^{n}:\left|z_{v}-a_{v}\right|<r_{v}, v=1, \ldots, n\right\}

\end{gathered}

\]



комплексного пространства $\mathbb{C}^{n}, n \geqslant 1$, являютцаяся топологич. произведением $n$ кругов,



\[

\begin{gathered}

\Delta=\Delta_{1} \times \ldots \times \Delta_{n}, \\

\Delta_{v}=\left\{z_{v} \in \mathbb{C}:\left|z_{v}-a_{v}\right|<r_{v}\right\}, v=1, \ldots, n .

\end{gathered}

\]



Точка $a=\left(a_{1}, \ldots, a_{n}\right) \in \mathbb{C}^{n}$ - центр поликруга $\Delta$, $r=\left(r_{1}, \ldots, r_{n}\right), r_{v}>0, v=1, \ldots, n,-$ его м у л ь т ир а д и у с. При $a=0, r=(1, \ldots, 1)$ получается е д иничны й полик р уг. Остов ом поликруг а $\Delta$ наз. часть



\[

T=T(a, r)=\left\{z \in \mathbb{C}^{n}:\left|z_{v}-a_{v}\right|=r_{v}, v=1, \ldots, n\right\}

\]



его полной топологич. границы $\partial \Delta$. П. есть полная кратно круговая область.



Естественным обобщением понятия П. является п оли о бласть (поликруговая область, обобше нный поли ци линд р) $D=D_{1} \times$ (. . $\times D_{n}$, являющаяся топологич. произведением

(вообе говоря, многосвязных ) областей $D \vee \subset \mathbb{C}, v=1$, ..., n. Граница $\Gamma=\partial D$ полиобласти $D$ состоит из $n$ множеств размерності $2 n-1$ :



$\Gamma_{v}=\left\{z \in \mathbb{C}^{n}: z_{v} \in \partial D_{v}, z_{\mu} \in \bar{D}_{\mu}, \mu \neq v\right\}, v=1, \ldots, m$, общая часть к-рых есть $n$-мерный остов полио б л а с т и $D$ :



$T=\partial D_{1} \times \ldots \times \partial D_{n}=\left\{z \in \mathbb{C}^{n}: z_{v} \in \partial D_{v}, v=1, \ldots, n\right\}$.



Е. Д. Соломенцев.



ПОЛИЛИНЕИНАЯ АЛГЕБРА - часть алгебры, изучающая полилинейные отображения модулей (в частности, векторных пространств). Первыми разделами П. а. явились теория билипейных борм и квадратичных форм, теория определителей и развивающее ее исчисление Грассмана (см. Внешняя алгебра). Основную роль в П. а. играют понятия тензорного произведения, тензора, полилинейной формы. Приложения П. а. к геометрии и анализу связаны главным образом с тензорныл исчислением и дифберенциальными бормами. А.Л. Онищик.



ПОЛИЛИНЕИНАЯ ФОРМА, $n-л и$ и е й н а я Ф о р м а, на унитарном $A$-модуле $E$ - полилинейное отображение $E^{n} \rightarrow A$ (здесь $A$ - ассоциативно-коммутативное кольцо с единицей). П. ф. наз. также по лили не й но й функци е й (n-линей ной ф ун к ц и е й). Поскольку П. Ф.- частный случай полилинейных отображений, можно говорить о симметрических, кососимметрических, знакопеременных, симметризованных и кососимметризованных П. Ф. Напр., определитель квадратной матрицы порядка $n$ над $A$ это кососимметризованная (и тем самым знакопеременная) $n$-линейная форма на $A^{n}$. $n$-линейные формы на $E$ образуют $A$-модуль $L_{n}(E, A)$, естественно изоморфный модулю $(\underbrace{n} E)^{*}$ всех линейных форм на $\bigotimes^{n} E$. В случае $n=2(n=3)$ говорят о билинейных бормах (трилинейных формах).



$n$-линейные формы на $E$ тесно связаны с $n$ раз ковариантными тензорами, т. е. элементами модуля $T^{n}\left(E^{*}\right)=\otimes^{n} E^{*}$. Точнее, имеется линейное отображение



$\quad \gamma_{n}: T^{n}\left(E^{*}\right) \longrightarrow L_{n}(E, A)$



такое, что



\[

\gamma_{n}: T^{n}\left(E^{*}\right) \longrightarrow L_{n}(E, A)

\]



\[

\gamma_{n}\left(u_{1} \otimes \ldots \otimes u_{n}\right)\left(x_{1}, \ldots, x_{n}\right)=u_{1}\left(x_{1}\right) \ldots u_{n}\left(x_{n}\right)

\]



для любых $u_{i} \in E^{*}, x_{i} \in E$. Если модуль $E$ свободен, то $\gamma$ инъективно, а если $E$ к тому же конечно порожден, то и биективно. В частности, $n$-линейные формы на конечномерном векторном пространстве над полем отождествляются с $n$ раз ковариантными тензорами.



Для любых форм $u \in L_{n}(E, A), v \in L_{m}(E, A)$ определяется их тензорное произведение $u \otimes v \in L_{n+m}(E, A)$ формулой



$u \otimes v\left(x_{1}, \ldots, x_{n+m}\right)=u\left(x_{1}, \ldots, x_{n}\right) v\left(x_{n+1}, \ldots, x_{n+m}\right)$. Для симметризованных П. ф. определено также симметрич. произведение



\[

\left(\sigma_{n} u\right) \vee\left(\sigma_{m} v\right)=\sigma_{n+m}(u \otimes v)

\]



а для кососимметризованных П. ф.- внешнее произведение



\[

\left(\alpha_{n} u\right) \wedge\left(\alpha_{m} v\right)=\alpha_{n+m}(u \otimes v) .

\]



Эти операции распространяются на модуль $L_{*}(E, A)=$ $=\bigoplus_{n=0}^{\infty} L(E, A)$, где $L_{0}(E, A)=A, L_{1}(E, A)=E^{*}$, модуль симметризованных форм $L_{\sigma}(E, A)=\bigoplus_{n=0}^{\infty} \sigma_{n} L_{n}(E$, $A)$ и модуль кососимметризованных форм $L_{\alpha}(E, A)=$ $=\bigoplus_{i=0}^{\infty} \alpha_{n} L_{n}(E, A)$ соответственно, превращая их в ассоциативные алгебры с единицами. Если $E-$ конечно порожденный свободный модуль, то отображения $\gamma_{n}$ определяют изоморфизм тензорной алгебры





\end{document}